\subsection{性能指標から有限ホライズン最適制御へ}

極配置では、閉ループ極を指定して応答の速さと減衰を調整した。しかし、質量 $m=1\,\mathrm{kg}$ の台車を目標位置で止める記録を評価すると、速さだけでは設計を決められないことが明確になる。位置偏差 $e$ [m] と速度偏差 $\dot{e}$ [m/s] を状態とし、入力 $u$ [N] を台車に加える水平力とする。強い入力を使えば整定は速くなるが、力の上限 $u_{\max}=2.0\,\mathrm{N}$ に近づき、RMS 入力も大きくなる。反対に入力を節約すれば余裕は増えるが、終端時刻で位置偏差が残りやすい。

\begin{table}[tb]
  \centering
  \small
  \caption{台車位置制御ログから得た性能指標の例。残留誤差、ピーク入力、RMS 入力、制約余裕は同時に改善しないため、単一の極位置だけでは設計判断を閉じられない。}
  \label{tab:flow-four-performance-indices}
  \begin{tabular}{llllll}
    \hline
    候補 & 整定時間 & 終端残差 & ピーク入力 & RMS 入力 & 入力余裕 \\
    \hline
    積極的 & $0.75\,\mathrm{s}$ & $0.004\,\mathrm{m}$ & $2.10\,\mathrm{N}$ & $1.25\,\mathrm{N}$ & $-0.10\,\mathrm{N}$ \\
    中庸 & $1.05\,\mathrm{s}$ & $0.006\,\mathrm{m}$ & $1.55\,\mathrm{N}$ & $0.82\,\mathrm{N}$ & $0.45\,\mathrm{N}$ \\
    穏やか & $1.85\,\mathrm{s}$ & $0.017\,\mathrm{m}$ & $0.95\,\mathrm{N}$ & $0.46\,\mathrm{N}$ & $1.05\,\mathrm{N}$ \\
    \hline
  \end{tabular}
\end{table}

\wtab{flow-four-performance-indices}では、積極的な設計は整定時間と終端残差に優れるが、ピーク入力が上限を越え、入力余裕 $M_u=u_{\max}-u_{\mathrm{peak}}$ が負になっている。穏やかな設計は制約余裕を大きく残す一方で、終端残差が大きい。RMS 入力
\begin{equation}
  u_{\mathrm{rms}}
  =
  \left(
    \frac{1}{T}
    \int_0^T u(t)^2\,\dd t
  \right)^{1/2}
\end{equation}
は、電流入力なら発熱、力入力なら入力努力の二乗積分に対応する負担を表す。このように、秒、メートル、ニュートンをそのまま足しても意味はないので、設計で比較するには各指標を要求値や定格値で割って無次元化する。

たとえば、要求整定時間 $T_{\mathrm{req}}$、許容終端誤差 $e_{\mathrm{tol}}$、入力上限 $u_{\max}$、連続定格入力 $u_{\mathrm{rated}}$、要求余裕 $M_{\mathrm{req}}$ を用いて
\begin{align}
  \widehat{T}_s&=\frac{T_s}{T_{\mathrm{req}}},
  &
  \widehat{e}_T&=\frac{\abs{e(T)}}{e_{\mathrm{tol}}},
  &
  \widehat{u}_{\mathrm{peak}}&=\frac{u_{\mathrm{peak}}}{u_{\max}},\\
  \widehat{u}_{\mathrm{rms}}&=\frac{u_{\mathrm{rms}}}{u_{\mathrm{rated}}},
  &
  \widehat{M}_u&=
  \frac{\max(0,M_{\mathrm{req}}-M_u)}{u_{\max}}
\end{align}
と置く。$\widehat{M}_u$ は余裕そのものではなく、必要な余裕を下回った分を罰する量である。これらはすべて無次元なので、重み $w_s,w_e,w_p,w_r,w_m\ge0$ を使って
\begin{equation}
  J_{\mathrm{log}}
  =
  w_s\widehat{T}_s^2
  +w_e\widehat{e}_T^2
  +w_p\widehat{u}_{\mathrm{peak}}^2
  +w_r\widehat{u}_{\mathrm{rms}}^2
  +w_m\widehat{M}_u^2
  \label{eq:normalized-log-cost}
\end{equation}
というスカラー評価を構成できる。重みは単位を消すための係数ではなく、すでに無次元化された指標の優先順位を表す係数である。$w_p$ や $w_m$ を大きくすれば入力上限への接近を強く避け、$w_s$ や $w_e$ を大きくすれば速い収束と残差低減を優先する。

ログ全体の評価を、時刻 $k$ の入力を選ぶ局所的な最適化へ置き換える。サンプリング周期を $h$ とし、台車を短い時間では二重積分器
\begin{align}
  e_{k+1}&=e_k+h\dot{e}_k+\frac{h^2}{2m}u_k,\\
  \dot{e}_{k+1}&=\dot{e}_k+\frac{h}{m}u_k
  \label{eq:cart-double-integrator-step}
\end{align}
で近似する。位置、速度、入力の代表スケールを $e_s$ [m]、$v_s$ [m/s]、$u_s$ [N] とし、
\begin{equation}
  z_k=
  \begin{bmatrix}
    e_k/e_s\\
    \dot{e}_k/v_s
  \end{bmatrix},
  \qquad
  v_k=\frac{u_k}{u_s}
\end{equation}
と無次元化する。すると
\begin{equation}
  z_{k+1}=\bar{A}z_k+\bar{B}v_k
\end{equation}
と書ける。ここで
\begin{equation}
  \bar{A}=
  \begin{bmatrix}
    1 & hv_s/e_s\\
    0 & 1
  \end{bmatrix},
  \qquad
  \bar{B}=
  \begin{bmatrix}
    h^2u_s/(2me_s)\\
    hu_s/(mv_s)
  \end{bmatrix}
\end{equation}
であり、$\bar{A}$ と $\bar{B}$ は無次元である。次の 1 サンプルだけを見て、次時刻の位置・速度偏差と入力努力を
\begin{equation}
  J_1(v_k)
  =
  z_{k+1}^\T W_x z_{k+1}
  +w_u v_k^2,
  \qquad
  W_x=W_x^\T\ge0,\quad w_u>0
  \label{eq:one-step-cost}
\end{equation}
で評価する。これはまだ LQR ではなく、表のような複数指標を、1 サンプル先の残差と入力努力へ落とした最小の有限ホライズン問題である。

\weq{one-step-cost} に $z_{k+1}=\bar{A}z_k+\bar{B}v_k$ を代入すると
\begin{align}
  J_1(v_k)
  &=
  (\bar{A}z_k+\bar{B}v_k)^\T
  W_x
  (\bar{A}z_k+\bar{B}v_k)
  +w_u v_k^2.
\end{align}
入力はスカラーなので、微分して
\begin{align}
  \frac{\partial J_1}{\partial v_k}
  &=
  2\bar{B}^\T W_x(\bar{A}z_k+\bar{B}v_k)
  +2w_u v_k\\
  &=
  2(\bar{B}^\T W_x\bar{B}+w_u)v_k
  +2\bar{B}^\T W_x\bar{A}z_k
\end{align}
を得る。最小化条件 $\partial J_1/\partial v_k=0$ から
\begin{equation}
  v_k^*
  =
  -(\bar{B}^\T W_x\bar{B}+w_u)^{-1}
  \bar{B}^\T W_x\bar{A}z_k
  \label{eq:one-step-optimal-input}
\end{equation}
である。分母に $w_u$ が入るため、入力を重く罰すると $v_k^*$ の大きさは小さくなる。たとえば $h=0.10\,\mathrm{s}$、$e_s=0.10\,\mathrm{m}$、$v_s=0.50\,\mathrm{m/s}$、$u_s=2.0\,\mathrm{N}$、$W_x=\diag(1,0.2)$、$z_k=[1\ 0]^\T$ では、$w_u=0.05$ なら $v_k^*=-1.09$ となり上限相当の $|v_k|=1$ を越える。$w_u=0.5$ では $v_k^*=-0.18$、$w_u=5$ では $v_k^*=-0.020$ であり、入力余裕は増えるが 1 サンプル後の残差は残りやすい。この数値は、重みが単なる記号ではなく、制約余裕と残留誤差の交換条件を直接動かすことを示している。

1 サンプルだけでは、少し先で速度を落として止まる計画を表しにくい。そこで、$N$ ステップ先までの入力列
\begin{equation}
  V_k=
  \begin{bmatrix}
    v_k\\
    v_{k+1}\\
    \vdots\\
    v_{k+N-1}
  \end{bmatrix}
\end{equation}
をまとめて選ぶ。予測状態を
\begin{equation}
  Z_k=
  \begin{bmatrix}
    z_{k+1}\\
    z_{k+2}\\
    \vdots\\
    z_{k+N}
  \end{bmatrix}
\end{equation}
と積むと、線形予測は
\begin{equation}
  Z_k=\mathcal{A}z_k+\mathcal{B}V_k
  \label{eq:stacked-prediction}
\end{equation}
で表される。行列は
\begin{equation}
  \mathcal{A}=
  \begin{bmatrix}
    \bar{A}\\
    \bar{A}^2\\
    \vdots\\
    \bar{A}^N
  \end{bmatrix},
  \qquad
  \mathcal{B}=
  \begin{bmatrix}
    \bar{B} & 0 & \cdots & 0\\
    \bar{A}\bar{B} & \bar{B} & \cdots & 0\\
    \vdots & \vdots & \ddots & \vdots\\
    \bar{A}^{N-1}\bar{B} & \bar{A}^{N-2}\bar{B} & \cdots & \bar{B}
  \end{bmatrix}
\end{equation}
である。途中の残差を $W_x$、終端残差を $S$、入力列を $w_u I$ で罰すれば
\begin{equation}
  J_N(V_k)
  =
  Z_k^\T Q_N Z_k+V_k^\T R_N V_k,
  \qquad
  Q_N=\diag(W_x,\ldots,W_x,S),
  \quad
  R_N=w_u I
\end{equation}
となる。代入して平方形式として微分すると
\begin{equation}
  V_k^*
  =
  -(\mathcal{B}^\T Q_N\mathcal{B}+R_N)^{-1}
  \mathcal{B}^\T Q_N\mathcal{A}z_k
  \label{eq:stacked-finite-horizon-solution}
\end{equation}
を得る。これが制約なしの多段有限ホライズン最適入力列である。$N$ を長くすると、入力を今だけでなく将来の停止準備にも配分できる。$R_N$ を大きくすれば列全体の入力努力が下がり、$S$ を大きくすれば終端時刻の残差を強く嫌う。この有限個のベクトル最適化を、時間を連続にし、非線形な状態方程式や一般のコストへ広げると、変分法、Pontryagin の最大原理、動的計画法が必要になる。二次コストと線形ダイナミクスに戻した特殊な場合が、後で導く LQR とリカッチ方程式である。

\subsection{変分法と最適性条件}

前段の 1 ステップ問題と多段予測問題では、入力は有限個の数 $v_k,\ldots,v_{k+N-1}$ だった。連続時間でホライズン全体を扱うと、選ぶ対象は入力列ではなく入力関数 $u(\cdot)$ になる。したがって評価関数は、数 $u$ の関数ではなく、関数 $u(\cdot)$ を受け取って数を返す汎関数である。典型的な有限時間問題は
\begin{align}
  J[u]
  &=
  \phi(x(T))
  +\int_0^T \ell(x(t),u(t),t)\,\dd t,\\
  \dot{x}(t)&=f(x(t),u(t),t),
  \qquad x(0)=x_0
  \label{eq:oc-general-problem}
\end{align}
で表される。$\ell$ は途中で支払うランニングコスト、$\phi$ は終端コストである。状態 $x(t)\in\R^n$ と入力 $u(t)\in\R^m$ は時刻に依存する関数であり、制御入力を少し変えると、状態軌道全体も変わる。この「関数を少し変えたときの評価関数の一次の変化」を調べる操作を変分という。

許される小さな入力変化を $\eta(t)$ とし、
\begin{equation}
  u_\varepsilon(t)=u(t)+\varepsilon \eta(t)
\end{equation}
と置く。対応する状態を $x_\varepsilon(t)$ と書くと、第一変分は
\begin{equation}
  \delta J
  =
  \left.\frac{\dd}{\dd\varepsilon}J[u_\varepsilon]\right|_{\varepsilon=0}
  \label{eq:first-variation}
\end{equation}
で定義される。通常の一変数関数で最小値の候補が $F'(z)=0$ を満たすように、制約のない十分滑らかな最適入力では、すべての許される $\eta$ に対して $\delta J=0$ が必要条件になる。この条件は必要条件であって、凸性や正定性がなければ大域最小性までは保証しない。

変分の基本形は、入力制御を含まない曲線最適化で確認できる。軌道 $x(t)$ そのものを選び、
\begin{equation}
  J[x]=\int_0^T L(x(t),\dot{x}(t),t)\,\dd t
\end{equation}
を最小にする問題を考える。端点 $x(0)$、$x(T)$ を固定し、$x_\varepsilon=x+\varepsilon \xi$、$\xi(0)=\xi(T)=0$ と置くと
\begin{align}
  \delta J
  &=
  \int_0^T
  \left(
    \frac{\partial L}{\partial x}\xi
    +\frac{\partial L}{\partial \dot{x}}\dot{\xi}
  \right)\,\dd t\\
  &=
  \int_0^T
  \left(
    \frac{\partial L}{\partial x}
    -\frac{\dd}{\dd t}\frac{\partial L}{\partial \dot{x}}
  \right)\xi\,\dd t
  +
  \left[
    \frac{\partial L}{\partial \dot{x}}\xi
  \right]_0^T.
\end{align}
端点固定により境界項は 0 である。任意の $\xi$ に対して $\delta J=0$ となるには、括弧内が各時刻で 0 でなければならない。よって Euler--Lagrange 方程式
\begin{equation}
  \frac{\partial L}{\partial x}
  -\frac{\dd}{\dd t}
  \frac{\partial L}{\partial \dot{x}}
  =0
  \label{eq:euler-lagrange-entry}
\end{equation}
を得る。これは「最適な曲線では、許される微小な曲線変化に対して評価関数の一次変化が消える」という条件である。力学で Euler--Lagrange 方程式が現れるのも同じ変分の構造による。

制御問題では、軌道 $x$ は自由に選べず、状態方程式 $\dot{x}=f(x,u,t)$ に従う。この動的制約を扱うため、随伴変数 $\lambda(t)\in\R^n$ を導入し、拡張された汎関数
\begin{equation}
  \bar{J}
  =
  \phi(x(T))
  +\int_0^T
  \left\{
    \ell(x,u,t)
    +\lambda^\T\bigl(f(x,u,t)-\dot{x}\bigr)
  \right\}\,\dd t
\end{equation}
を考える。最適軌道では $f-\dot{x}=0$ なので、$\bar{J}$ と $J$ の値は一致する。$x$ と $u$ の変分を取り、$\lambda^\T\dot{x}$ の項を部分積分して一次の項を集めると、Pontryagin の Hamiltonian
\begin{equation}
  \mathcal{H}(x,u,\lambda,t)
  =
  \ell(x,u,t)+\lambda^\T f(x,u,t)
  \label{eq:pontryagin-hamiltonian}
\end{equation}
により、内部点の必要条件は
\begin{align}
  \dot{x}&=\frac{\partial \mathcal{H}}{\partial \lambda}
  =f(x,u,t),
  \label{eq:pmp-state}\\
  -\dot{\lambda}&=\frac{\partial \mathcal{H}}{\partial x},
  \label{eq:pmp-costate}\\
  0&=\frac{\partial \mathcal{H}}{\partial u},
  \label{eq:pmp-stationarity}\\
  \lambda(T)&=\frac{\partial \phi}{\partial x}(x(T)).
  \label{eq:pmp-terminal}
\end{align}
となる。$\lambda$ は状態を少し変えたときに将来の最小コストがどれだけ変わるかを表す感度であり、随伴変数または共状態と呼ばれる。状態方程式は初期条件 $x(0)=x_0$ から前向きに進み、随伴方程式は終端条件 \weq{pmp-terminal} から後ろ向きに決まる。したがって、最適制御問題は一般に二点境界値問題になる。入力制約がない滑らかな場合は \weq{pmp-stationarity} を使うが、入力に上下限や不等式制約がある場合は、Hamiltonian を許容集合上で最小にする条件や KKT 条件へ置き換わる。

同じ問題の別表現が動的計画法である。時刻 $t$、状態 $x$ から終端までに達成できる最小残りコストを値関数
\begin{equation}
  V(t,x)=
  \min_{u(\cdot)}
  \left\{
    \phi(x(T))
    +\int_t^T \ell(x(\tau),u(\tau),\tau)\,\dd \tau
  \right\}
  \label{eq:value-function-entry}
\end{equation}
と定義する。短い時間幅 $h>0$ だけ先へ進めると、最適性の原理より
\begin{equation}
  V(t,x)=
  \min_{u(\cdot)}
  \left\{
    \int_t^{t+h}\ell(x(\tau),u(\tau),\tau)\,\dd \tau
    +V(t+h,x(t+h))
  \right\}
  \label{eq:dynamic-programming-principle}
\end{equation}
が成り立つ。右辺は、最初の短い区間で支払うコストと、その後の最小残りコストの和である。$h$ を十分小さくし、$x(t+h)=x+h f(x,u,t)+o(h)$ を代入して Taylor 展開すると
\begin{align}
  0=
  \min_u
  \left\{
    \ell(x,u,t)
    +\frac{\partial V}{\partial t}(t,x)
    +\frac{\partial V}{\partial x}(t,x)^\T f(x,u,t)
  \right\}
  \label{eq:hjb-entry}
\end{align}
を得る。これが Hamilton--Jacobi--Bellman 方程式であり、終端条件は $V(T,x)=\phi(x)$ である。Pontryagin の条件が一つの最適軌道に沿った必要条件を与えるのに対し、HJB 方程式は全状態に対する値関数を求める方程式である。滑らかな値関数が得られる場合、両者は
\begin{equation}
  \lambda(t)=
  \frac{\partial V}{\partial x}(t,x(t))
  \label{eq:costate-value-gradient}
\end{equation}
で結びつく。

LQR では、この一般形が閉じた行列方程式に落ちる。線形系 $f=Ax+Bu$、二次コスト
\begin{equation}
  \ell=x^\T Qx+u^\T Ru,
  \qquad
  \phi=x^\T Sx
\end{equation}
を用いると、Hamiltonian は
\begin{equation}
  \mathcal{H}
  =
  x^\T Qx+u^\T Ru+\lambda^\T(Ax+Bu)
\end{equation}
であり、停留条件から
\begin{equation}
  0=2Ru+B^\T\lambda,
  \qquad
  u=-\frac{1}{2}R^{-1}B^\T\lambda
\end{equation}
を得る。一方、HJB で値関数を $V(t,x)=x^\T P(t)x$ と置けば
\begin{equation}
  \lambda=\partial_x V=2P(t)x
\end{equation}
なので、最適入力は
\begin{equation}
  u=-R^{-1}B^\T P(t)x
\end{equation}
となる。さらに HJB 方程式の二次形式の係数を 0 と置くことでリカッチ微分方程式が得られ、無限ホライズンではその定常解が代数リカッチ方程式になる。したがって、変分、Pontryagin の条件、動的計画法、HJB は別々の公式ではなく、LQR のリカッチ方程式へ至る同じ最適性の構造を異なる表現で見ている。

最後に、一次元の積分器でこの接続を確認する。対象を
\begin{equation}
  \dot{x}=u,
  \qquad x(0)=x_0
\end{equation}
とし、評価関数を
\begin{equation}
  J=sx(T)^2+\int_0^T r u(t)^2\,\dd t,
  \qquad s\ge0,\quad r>0
  \label{eq:scalar-integrator-cost}
\end{equation}
とする。Hamiltonian は
\begin{equation}
  \mathcal{H}=r u^2+\lambda u
\end{equation}
である。最適性条件は
\begin{align}
  \dot{x}&=u,\\
  -\dot{\lambda}&=0,\\
  0&=2ru+\lambda,\\
  \lambda(T)&=2s x(T)
\end{align}
である。したがって $\lambda$ は定数で、$u=-\lambda/(2r)$ も定数になる。終端条件に $x(T)=x_0+Tu$ を代入すると
\begin{align}
  \lambda
  &=2s\left(x_0-\frac{T}{2r}\lambda\right),\\
  \lambda\left(1+\frac{sT}{r}\right)&=2s x_0,\\
  u^*&=-\frac{s}{r+sT}x_0
\end{align}
を得る。同じ結果を HJB から導く。$V(t,x)=P(t)x^2$ と置いたとき
\begin{equation}
  0=\min_u\{r u^2+\dot{P}x^2+2P xu\}
\end{equation}
であり、$u^*=-(P/r)x$、さらに
\begin{equation}
  \dot{P}=\frac{P^2}{r},
  \qquad
  P(T)=s
\end{equation}
となる。このスカラーのリカッチ方程式を解くと
\begin{equation}
  P(t)=\frac{sr}{r+s(T-t)}
  \label{eq:scalar-riccati-solution}
\end{equation}
である。したがって $t=0$ では
\begin{equation}
  u^*(0)=-\frac{P(0)}{r}x_0
  =-\frac{s}{r+sT}x_0
\end{equation}
となり、Pontryagin の条件から得た入力と一致する。この小さな例では、随伴変数、値関数の勾配、リカッチ方程式が同じ最適性条件の別表現であることが直接確認できる。

\subsection{確率とノイズ}

ここまでの制御では、状態 $x$ が分かっているか、少なくともオブザーバで滑らかに推定できると考えた。台車の例では、位置センサから位置は読めても、速度は測れないことが多い。さらに、センサの分解能や電気的な揺らぎにより、測定値をそのまま真値とみなすと入力が細かく振れる。測定値 $y$ を
\begin{equation}
  y=Cx+v
\end{equation}
と書くと、$v$ は測定ノイズである。$y$ の単位は測定量の単位であり、位置センサならメートル、角度センサならラジアンである。平均が 0、分散が $\sigma^2$ のノイズなら
\begin{equation}
  E[v]=0,
  \qquad
  E[v^2]=\sigma^2
\end{equation}
である。多次元では共分散行列
\begin{equation}
  R=E[vv^\T]
\end{equation}
がノイズの大きさと相関を表す。$R$ の $(i,j)$ 成分は、測定 $y_i$ と $y_j$ の単位の積を持つ。

モデルにも不確かさがある。摩擦、坂道、外乱、離散化誤差をすべて正確には書けないため、離散時間モデルを
\begin{align}
  x_{k+1}&=A x_k+B u_k+w_k,\\
  y_k&=C x_k+v_k
\end{align}
と書き、$w_k$ の共分散を $Q$、$v_k$ の共分散を $R$ とする。この導入では記号を簡潔に $Q,R$ と書くが、後の詳細な導出では LQR の重みと区別するためにプロセスノイズ共分散を $Q_w$、測定ノイズ共分散を $R_v$ と書く。$Q$ の成分は状態成分の単位の積、$R$ の成分は測定成分の単位の積を持つ。$Q$ が大きいほどモデルを疑い、$R$ が大きいほどセンサを疑う。推定とは、この二つの信頼度を使って状態を更新することである。

\subsection{カルマンフィルタ}

まず1次元で、予測値と測定値の混ぜ方を確認する。事前推定値を $\hat{x}^-$、その分散を $P^-$、測定値を $y$、測定ノイズ分散を $R$ とする。$x$ が位置なら $P^-$ と $R$ は平方メートルの単位を持ち、値が大きいほどその情報を信じにくいことを表す。更新後の推定値を
\begin{equation}
  \hat{x}=\hat{x}^-+K(y-\hat{x}^-)
\end{equation}
と置くと、$K$ は測定をどれだけ信じるかを表す。誤差分散を最小にするゲインは
\begin{equation}
  K=\frac{P^-}{P^-+R}
\end{equation}
である。事前分散が大きいほど測定を信じ、測定ノイズが大きいほど事前値を信じる。

多次元のカルマンフィルタは、線形モデルに対する予測と更新で書ける。
\begin{align}
  \hat{x}_{k|k-1}&=A\hat{x}_{k-1|k-1}+Bu_{k-1},\\
  P_{k|k-1}&=AP_{k-1|k-1}A^\T+Q,\\
  K_k&=P_{k|k-1}C^\T(CP_{k|k-1}C^\T+R)^{-1},\\
  \hat{x}_{k|k}&=\hat{x}_{k|k-1}+K_k(y_k-C\hat{x}_{k|k-1}),\\
  P_{k|k}&=(I-K_kC)P_{k|k-1}.
\end{align}
ここで $P$ は推定誤差の共分散である。カルマンフィルタは、線形ガウス系では平均二乗誤差を最小にする推定器である。

LQR は状態が利用可能である場合の最適状態フィードバックであり、Kalman フィルタは線形確率モデルにおける状態推定器である。この接続はここでは前方参照にとどめ、後続の節で LQR と Kalman フィルタをそれぞれ動機づけ、式として導いた後に、推定値 $\hat{x}$ を状態フィードバックへ渡す構成を LQG として扱う。本節では、その準備として対象および推定器を線形として扱う。
